% Generated from validated Article Draft Result. Do not hand-edit; revise the structured draft instead.
\section{AI for Science — 4月、科学AIは「専門化」のフェーズへ}
\label{pkg:ai-science-april}
\noindent\textbf{GPT-Rosalind、SPUR、DPLM-Evo。4月に見えたのは完全自律研究室ではなく、科学workflow向けモデル、scientific image benchmark、protein evolution生成という専門化の三方向だった。}\autocite{src-078f0d540d88,src-7ae707d75c4b,src-f0bb84fcd319}

\subsection{科学領域では「汎用モデルの応用」から専門化へ}

4月のAI for Scienceを、すでにAIが研究室を自律運転した月として描くのは正確ではない。この時点で見えるのは、創薬・genomics・protein reasoningへscopeを絞ったfrontier model、scientific experimental imageを測るbenchmark、protein evolutionを生成問題として扱うdiffusion researchという三つの方向である。共通するのは「科学」を一つのbenchmark domainとして扱うのではなく、対象データとworkflowに合わせてmodel／evaluationを専門化する動きが強まったことだ。\autocite{src-078f0d540d88,src-7ae707d75c4b,src-f0bb84fcd319}


\subsection{GPT-Rosalind — scientific workflowを明示したfrontier model}

OpenAIは4月16日にGPT-Rosalindを公表し、drug discovery、genomics analysis、protein reasoning、scientific research workflowsを対象scopeとして説明した。この一次資料から確実に固定できるのはrelease chronologyとOpenAIが示した用途範囲である。個別taskでの性能、実験室での自律性、研究成果への一般的な寄与をここから独立事実として導くことはできない。4月号では、frontier reasoning modelが専門科学領域を明示的なproduct／research surfaceとして切り出したことに意味を置く。\autocite{src-078f0d540d88}


\subsection{SPUR — scientific imageを独立した評価対象にする}

SPURは4月30日の初回投稿で、scientific experimental imageのperception、understanding、reasoningを評価するbenchmarkとして提示された。一般的なVQAや自然画像理解だけでは、顕微鏡像、実験図、計測画像など科学固有の視覚情報をどこまで扱えるかを十分に測れないという問題意識がある。したがってSPURの重要性は特定modelの順位だけでなく、scientific image interpretationを独立したevaluation problemとして定義した点にある。保持locatorは5月26日のlater revisionであり、4月の初稿と後日更新を混同しない。\autocite{src-7ae707d75c4b}


\subsection{DPLM-Evo — protein evolutionを生成過程として捉える}

DPLM-Evoは4月30日の初回投稿で、denoising中にsubstitution、insertion、deletionを明示的に予測するevolutionary discrete diffusion frameworkを提案した。通常の固定長sequence generationとは異なり、進化に現れる編集操作そのものを生成過程へ組み込むという問題設定である。これはGPT-Rosalindのような汎用reasoning modelの科学適用とは別の研究線で、domain-specific generative modelingの例として位置付ける。保持locatorは6月3日のlater revisionであるため、後日の実験記述を4月へ遡及しない。\autocite{src-f0bb84fcd319}


\subsection{model、benchmark、generative mechanismの三方向}

GPT-Rosalind、SPUR、DPLM-Evoを一つの「科学AI」ランキングへ並べることには意味がない。GPT-Rosalindはscientific workflowをscopeに持つfrontier reasoning model、SPURはscientific image理解を測るbenchmark、DPLM-Evoはprotein sequence evolutionへ特化したgenerative frameworkである。むしろ三者の違いは、AI for Scienceがmodel capabilityだけでなく、evaluation targetとdomain mechanismへ分解されていく様子を示す。4月の「競争軸の分化」という号全体の主題が、科学領域にも現れている。\autocite{src-078f0d540d88,src-7ae707d75c4b,src-f0bb84fcd319}


\subsection{6月の実験workflowへつながるが、同じ段階ではない}

後続の6月号ではnear-autonomous chemistry workflowが独立Featureとなり、humanがproposal selection、lab operation、validationを担う境界まで含めて実世界workflowへの接続を扱う。4月の資料はそこまでを示していない。4月はmodel／benchmark／domain-specific generationが専門化した段階、6月はそれらの能力を長時間の実験workflowへ接続する段階、と分けることで進展の意味が明確になる。\autocite{src-078f0d540d88,src-7ae707d75c4b,src-f0bb84fcd319}


\begin{claimboundary}
Claim Boundary: GPT-Rosalindの一次資料はrelease chronologyとOpenAIが述べたscientific scopeを確認するものだが、capability・performance・実験上の有効性を独立再現したものではない。\autocite{src-078f0d540d88}
\end{claimboundary}

\begin{claimboundary}
Chronology Boundary: later arXiv revisionを保持する場合、April初回投稿の内容と後日改訂を分離する。SPURは2026年5月26日、DPLM-Evoは2026年6月3日に保持locatorが更新されており、post-Aprilの追加内容を4月時点へ帰属させない。\autocite{src-7ae707d75c4b,src-f0bb84fcd319}
\end{claimboundary}
